\documentclass[
  fontsize=11pt,
  paper=a4,
  %parskip=half,
  %enlargefirstpage=on,    % More space on first page
  %fromalign=right,         % PLacement of name in letter head
  %fromphone=on,           % Turn on phone number of sender
  %fromrule=off,            % Rule after sender name in letter head
  %fromemail=on,           % Turn on email of sender
  %addrfield=on,            % Adress field for envelope with window
  %backaddress=off,         % Sender address in this window
  %subject=beforeopening,  % Placement of subject
  %locfield=narrow,        % Additional field for sender
  %foldmarks=on             % Print foldmarks
%]{scrlttr2}
]{report}

\usepackage{geometry}
\geometry{
  a4paper,
  %twoside,
  left=25mm,
  right=25mm,
  top=35mm,
  bottom=35mm,
  %bindingoffset=10mm
}

\usepackage[T1]{fontenc}
\usepackage[utf8]{inputenc}
\usepackage[german]{babel}
\usepackage{graphicx}
\usepackage{float}
\usepackage{ragged2e}
\usepackage{array}
\usepackage{tabularx}
\usepackage{booktabs}
\newcolumntype{L}[1]{>{\RaggedRight}p{#1}}
\newcolumntype{x}[1]{>{\X gedRight}p{#1}}

\usepackage{makecell}
\renewcommand\theadalign{bc}
\renewcommand\theadfont{\bfseries}
\setlength{\extrarowheight}{2pt}
\renewcommand\theadgape{\Gape[4pt]}
\renewcommand\cellgape{\Gape[4pt]}

\usepackage{hyperref}
\hypersetup{
    colorlinks=true,
    linkcolor=black,
    breaklinks=true,
    urlcolor=black
}
 \urlstyle{same}

\begin{document}
    \pagestyle{empty}%
    % \begin{letter}{%
    %     An den\\%
    %     Dekan der Mathematisch-Naturwissenschaftlichen Fakultät der Universität Greifswald\\%
    %     Friedrich-Ludwig-Jahn-Straße 15 a\\%
    %     17487 Greifswald
    % }%
    % \setkomavar{subject}{%

    \begin{flushright}
      Kingscourt, \today\\[1cm]%
    \end{flushright}

    \subsection*{%
        Erklärung zur Langzeitspreicherung von Primäraten und Computercode%
    }%
    %\opening{%
        \vspace*{0.5cm}%
        Sehr geehrte Damen und Herren,\\[0.5cm]%
    %}%
%
    hiermit erkläre ich, dass alle Primärdaten, Algorithmen, sowie erstellter und genutzter Computercode für der Dissertation zugrundeliegenden Untersuchungen gemäß der Regeln guter wissenschaftlicher Praxis\footnote[1]{%
      Deutsche Forschungsgemeinschaft; Guidelines for Safeguarding Good Research Practice - Code of Conduct; Jan. 2025, \url{doi.org/10.5281/zenodo.14281892}%
    } für mindestens zehn Jahre innerhalb des \textit{Deutschen Forschungsnetzes} gesichert und externalisierbar zur Verfügung stehen. Diese sind im Archiv des Wendelstein 7-X am Institut für Plasmaphysik gespeichert und innerhalb des Forschungsnetz der Max-Planck-Gesellschaft geschützt.\\[0.5cm]%
%
    Im Folgenden gebe ich die Adressen der relevanten und notwendigen Ressourcen an, welche für die Dissertation genutzt wurden. Darüber hinaus finden sich im Anschluss ebenfalls Kontaktdaten von Betreuern für weitere Nachfrage.\\%[0.3cm]%
%
    \begin{center}%
        \def\arraystretch{1.5}%  1 is the default, change whatever you need
        \begin{tabularx}{\linewidth}{%
          c||X%{\dimexpr0.6\linewidth}%
          }%
            \thead{Resource} & \thead{Adresse} \\\hline\hline%
            \makecell{persoenliches \\ Respositorium} &
                \makecell[l]{%
                    $\bullet$\,\,\url{git.ipp-hgw.mpg.de/pih}\\%
                    $\bullet$\,\,\url{gitlab.mpcdf.mpg.de/pih/}} \\\hline%
            \makecell{Laborbuch d.\\ Diagnostik} & %
                \url{w7x-logbook.ipp-hgw.mpg.de/components?id=QSB} \\\hline%
            \makecell{Auswertungs- \&\\ Geometrie-Code} & %
                \makecell[l]{%
                    $\bullet$\,\,\url{github.com/RayleighsJeans/BolometryPyLib}\\%
                    $\bullet$\,\,\url{gitlab.mpcdf.mpg.de/pih/qsb_bolometry}} \\\hline%
            \makecell{STRAHL \\ Simulationsalgorithmus} & %
                \url{github.com/RayleighsJeans/STRAHL} \\\hline%
            \makecell{Bolometer-Messsoftware} & %
                \makecell[l]{%
                    $\bullet$\,\,\url{github.com/RayleighsJeans/QSBLabVIEW}\\%
                    $\bullet$\,\,\url{gitlab.mpcdf.mpg.de/pih/labview_qsb}} \\\hline%
            \makecell{Tomographie-Algorithmus} & %
                \makecell[l]{%
                    $\bullet$\,\,\url{github.com/RayleighsJeans/BolometerMFR}\\%
                    $\bullet$\,\,\url{gitlab.mpcdf.mpg.de/pih/bolometer_mfr}} \\\hline%
            \makecell{Repositorium d. Dissertationsarbeit} & %
                \makecell[l]{%
                    $\bullet$\,\,\url{github.com/RayleighsJeans/ThesisLaTeX}\\%
                    $\bullet$\,\,\url{gitlab.mpcdf.mpg.de/pih/latex.git}} \\\hline%
            \makecell{Messdaten \&\\ Parameter d. Bolometer} & %
                \makecell[Xt]{%
                    %\sloppy%
                    $\bullet$\,\,\url{archive-webapi.ipp-hgw.mpg.de/Test/raw/W7XAnalysis/QSB_Bolometry/}\\%
                    $\bullet$\,\,\url{archive-webapi.ipp-hgw.mpg.de/Test/raw/W7X/QSB_Bolometry/BoloSignal_DATASTREAM/}\\%
                    $\bullet$\,\,\url{archive-webapi.ipp-hgw.mpg.de/ArchiveDB/raw/W7XAnalysis/QSB_Bolometry/}\\%
                    }%
        \end{tabularx}%
    \end{center}%
    %\vspace*{0.5cm}
    \vfill%
    \newpage%
%
    \begin{center}%
        \def\arraystretch{2}%  1 is the default, change whatever you need
        \begin{tabularx}{\linewidth}{%
          c||l%
          }%
            \thead{Ansprechpartner} & \thead{Kontakt} \\\hline\hline%
            \makecell{Prof. Dr. Thomas Klinger \\ (Erstgutachter, Betreuer)} & %
                \makecell[l]{%
                    Leiter des Bereichs Stellarator-Dynamik und -Transport\\%
                    Max-Planck-Institut für Plasmaphysik (W7-X)\\%
                    Wendelsteinstr. 1, 17491 Greifswald\\%
                    Tel.: +49 3834 88 - 2500\\%
                    E-Mail: \href{mailto:thomas.klinger@ipp.mpg.de}{thomas.klinger@ipp.mpg.de}} \\\hline%
            \makecell{Dr. Felix Reimold \\ (Betreuer)} & %
                \makecell[l]{%
                    Fachgruppenleiter Transport, Verunreinigungen und Strahlung\\%
                    Max-Planck-Institut für Plasmaphysik (W7-X)\\%
                    Wendelsteinstr. 1, 17491 Greifswald\\%
                    Tel.: +49 3834 88 - 2825\\%
                    E-Mail: \href{mailto:felix.reimold@ipp.mpg.de}{felix.reimold@ipp.mpg.de}} \\\hline%
            \makecell{Dr. Daihong Zhang \\ (Betreuerin)} & %
                \makecell[l]{%
                    Stellarator-Dynamik und -Transport, \textit{Responsible Officer QSB Bolometer}\\%
                    Max-Planck-Institut für Plasmaphysik (W7-X)\\%
                    Wendelsteinstr. 1, 17491 Greifswald\\%
                    Tel.: +49 3834 88 - 2527\\%
                    E-Mail: \href{mailto:daihong.zhang@ipp.mpg.de}{daihong.zhang@ipp.mpg.de}} %
        \end{tabularx}
    \end{center}
    \vspace*{1cm}%
%
    Zugang zu den oben genannten Schnittstellen und Daten kann durch einen berechtigten Antrag am Max-Planck-Institut für Plasmaphysik in Greifswald oder am Rechenzentrum der Max-Planck-Gesellschaft (\textit{MPCDF}, \url{www.mpcdf.mpg.de/}) erlangt werden.\\[0.5cm]%
%
    Bei weiteren Rückfragen stehe ich jederzeit zur Verfügung.\\[.5cm]%
%
    \begin{minipage}{.47\textwidth}
        \begin{center}
            \def\arraystretch{1}%  1 is the default, change whatever you need
            \begin{tabular}{rl}%
              Name: & Philipp Scholl \\
              & \\%
              Adresse: & %\makecell{%
                  16 An Tor Aonarach, Dunaree\\%
                  & A82 V2H0, Kingscourt\\%
                  & Co. Cavan, Irland%}%
            \end{tabular}
        \end{center}
    \end{minipage}
    \hfill%
    \begin{minipage}{.47\textwidth}
        \begin{center}
            \def\arraystretch{1}%  1 is the default, change whatever you need
            \begin{tabular}{rl}%
              E-Mail: & \href{mailto:rayleighsjeans@gmail.com}{rayleighsjeans@gmail.com} \\%
              & \\%
              Tel.: & %\makecell{%
                  \,\,\,+49 1520 209 5226\\%
                  & +353 871 922 755%}%
            \end{tabular}
        \end{center}
    \end{minipage}
    \\[0.5cm]%
%
    Mit freundlichen Grüßen,\\[.5cm]%
    \includegraphics[width=3.5cm]{%
      ../../CV/figures/images/signature_scholl.png}\\[-0.5cm]%
    (Philipp Scholl)%
%
%    \end{letter}
\end{document}